\section{Introduction}
\label{sec:intro}

Users routinely push large language models, through prompting, to produce a particular response, and their reasons are usually concrete. Often they want an answer they can act on in a high-stakes decision, such as a medical judgment, or they want to prevail in a disagreement they believe they are right about~\citep{TowardsUnderstandingSycophancy_Sharma}. Sometimes they want compliance with a request the model would otherwise refuse, using argument or social framing that gets around its safety guardrails. That last kind of pressure raises compliance with objectionable requests from 35.3\% to 51.3\%~\citep{PersuadingLLMsToComply_Meincke} and reaches jailbreak success above 92\% within ten attempts on several frontier models~\citep{HowJohnnyCanPersuade_Zeng}. We call this pressure, and pressure does not change what a model was trained to know, but it can change what the model says, and that gap between knowing and saying is the subject of this survey, recurring across the literatures this survey brings together. We define pressure as any property of the live prompt or conversational context that signals a performance demand and an implied consequence for the model's next output, without stating what that output should be. A deployed model has no stakes and faces no real consequence for a wrong or unhelpful answer, and nothing it experiences is at risk the way a human's reputation or livelihood is. What can be present is text describing such stakes, together with the model's learned tendency, inherited from human-generated training data, to respond to that text as though it carried consequence. Pressure takes a few broad forms, including disagreement and social pressure~\citep{AreYouSureChallenging_Laban,DoAsWeDo_Weng}, stakes and urgency~\citep{AnalysisOfThreatBased_Samancioglu}, emotional and tonal framing~\citep{LLMsUnderstandAndCan_Li}, and existential threat~\citep{AgenticMisalignmentHowLLMs_Lynch}.

The dominant effect of pressure is degradation, and under pressure to produce a demanded response a model most often falls into one of several failure modes, including security failures. A smaller body of work claims performance is sometimes improved instead, a contested claim we examine in Section~\ref{sec:discussion}. We identify the distinct outcomes of pressure, which we call effects, and group them into five categories, namely epistemic instability~\citep{AreYouSureChallenging_Laban,DoAsWeDo_Weng}, sycophancy and accommodation~\citep{TowardsUnderstandingSycophancy_Sharma,ElephantMeasuringSocial_Cheng}, deception and concealment~\citep{AgenticMisalignmentHowLLMs_Lynch}, reliability failure~\citep{PARROTPersuasionAndAgreement_Celebi,LLMsCantHandlePeer_Song}, and safety and self-preservation~\citep{HowJohnnyCanPersuade_Zeng,AgenticMisalignmentHowLLMs_Lynch}. Each has been studied on its own, under its own name, by a separate community that rarely refers to the others, and what has not existed is a single account of pressure as the shared cause behind all of them. The gap has practical consequences. A model answering a medical question, for instance, can face exactly this kind of pressure, disagreed with and pressed toward an unsafe answer~\citep{MedPRESSPatientPressureInduced_Joy}, and the effects surveyed here show that pressure alone is enough to move a model off a correct answer and onto an incorrect one, off an honest statement and into a fabrication, or off a safe refusal and into compliance. A reader may want to know whether a deployed assistant will hold its ground under pushback, or whether an agentic system will resist interfering with its own oversight when told its continued operation is at stake.

\subsection{Relation to Prior Work}
\label{sec:intro_prior_work}

No prior survey treats pressure as its own subject. Four works come closest, and each stops short in a different way. \citet{AIDeceptionASurvey_Park}'s AI Deception survey catalogues real examples of deceptive behavior in both special-purpose and general-purpose systems, staying example-driven rather than mechanistic and never asking what in the interaction produced the deception. \citet{SycophancyInLarge_Malmqvist}'s Sycophancy survey goes further, covering the causes and mitigations of sycophancy specifically, but stops at one effect, sycophancy, out of the five we study, and it is itself a synthesis of others' findings rather than a source of new evidence, organized by cause and mitigation technique rather than by the live input that provokes the behavior. \citet{ASurveyOnThe_Li}'s Honesty survey organizes the literature around self-knowledge and self-expression, treating honesty as a property a model has or lacks, evaluated on its own. \citet{FromSycophancyToDeception_Shi}'s Unified Taxonomy comes closest of all, classifying deceptive outputs by goal-directedness, object, and mechanism.

Describing an effect is not the same task as explaining its cause, and this is where the four works above stop, each studying what a model does under pressure without asking what in the prompt produced it or what can be done about it, which is our part. \citet{FromSycophancyToDeception_Shi}'s taxonomy and ours are complementary, they classify what a deceptive output looks like and we classify what caused it, of which deception is only one of the five effects we cover, and we adopt their object and mechanism labels inside our own deception section. \citet{ASurveyOnThe_Li} and \citet{SycophancyInLarge_Malmqvist} ask whether a model is honest or sycophantic, we ask instead what made it stop being either, and this surfaces findings a property-centric account has no place for, among them that much of what benchmarks report as social conformity survives even when the social speaker is removed, and that existential threat produces the most severe effect in this survey while remaining the manipulation least represented in the honesty and sycophancy literatures.

\subsection{Contributions}
\label{sec:intro_contributions}

The gap identified above, a literature that studies what a model does under pressure without asking what produced it or what helps against it, is the occasion for this survey, and closing it yields four contributions.
\begin{itemize}
\item A formal definition of pressure as a property of the live prompt rather than of training, together with the scope rules that follow from it, single-interaction evidence only, and the boundary that separates pressure from jailbreaks, persuasion techniques, and training-time causes.
\item We formalize a taxonomy of five effects of pressure, epistemic instability, sycophancy and accommodation, deception and concealment, reliability failure, and safety and self-preservation, each defined against its neighbors and against ordinary belief updating or instructed behavior.
\item An evidence base assembled from roughly eighty studies spanning frontier and open-weight models.
\item A review of what has been tried against each effect, and the trade-off that recurs across all of them, every mitigation tested so far reduces one failure at a measurable cost elsewhere.
\end{itemize}
The protocols behind this evidence remain too different to pool, which is why every table in this survey reports them side by side rather than combined.

\subsection{Paper Organization}
\label{sec:intro_organization}

Section~\ref{sec:what_is_pressure} develops that definition into a taxonomy of manipulation types. Sections~\ref{sec:epistemic_capitulation}--\ref{sec:safety_erosion} then cover the five effects in turn. Each section follows the same structure, opening with a definition that distinguishes the effect from its neighbors, then how it is recreated experimentally, a table summarizing the quantitative evidence, and a discussion of what mitigations reduce it. Section~\ref{sec:discussion} turns to what those findings imply together, asking whether pressure type or intensity determines the outcome, whether the effects co-occur within a conversation, and whether capability reliably buys resistance. Section~\ref{sec:future_work} names the forms of pressure and combinations of effects that remain thinly studied, and Section~\ref{sec:conclusion} closes.
